\chapter*{Appendix A: Additional Experimental Results}
\addcontentsline{toc}{chapter}{Appendix A: Additional Experimental Results}
\begingroup
\justifying
\setlength{\parindent}{0pt}
\setstretch{2}
\setlength{\parskip}{0.5\baselineskip}

\section{Sample Predictions}

This section provides additional qualitative examples of predictions from the base model and the PPO fine-tuned model.

\begin{table}[h]
\centering
\caption{Sample Prediction Comparison}
\label{tab:sample_predictions}
\begin{tabular}{p{2cm}p{3cm}p{3cm}p{2cm}}
\toprule
\textbf{Sample ID} & \textbf{Ground Truth} & \textbf{Base Model} & \textbf{PPO Model} \\
\midrule
0 & omission & logical error, omission & logical error, omission \\
1 & logical error, omission & omission & logical error, omission \\
2 & omission & logical error, omission & logical error, omission \\
3 & none & logical error, omission & logical error, omission \\
\bottomrule
\end{tabular}
\end{table}

As shown in Table \ref{tab:sample_predictions}, the PPO model generally produces more accurate predictions, particularly for samples with the \textit{omission} label.

\section{Training Curves}

Figure \ref{fig:training_curves} shows the training curves for the PPO fine-tuning process, including the reward scores from the fluency and alignment discriminators.

\begin{figure}[h]
    \centering
    \includegraphics[width=0.6\textwidth]{assets/figures/training_curves.png}
    \caption{PPO Training Curves}
    \label{fig:training_curves}
\end{figure}

The training curves show a gradual increase in both fluency and alignment rewards over training, indicating that the model is learning to generate more fluent and accurate label sequences.

\endgroup
